\section{Evaluation Pipeline}

\subsection{Combining Sentiment (Acoustic) Output with the LLM Evaluation Pipeline}

The evaluation graph runs two independent signal sources per call: an LLM
reads the transcript against the rubric (compliance, quality, workflow,
escalation), and a separate acoustic sentiment model scores each
speaker-attributed sentence for valence, dominant emotion, and an
escalation score. These are combined through a deterministic \emph{fusion}
step rather than by feeding acoustic output back into the LLM prompt, so
that the acoustic signal's variance cannot silently perturb the text
model's judgment, and every fused score remains reproducible and
auditable after the fact.

Design principle: the fusion module is deterministic \emph{only where
agreement gives it something to be deterministic about}. Where the two
signal sources disagree, the disagreement is routed to a dedicated
arbitration node in the graph rather than resolved by a fixed rule --
disagreement is precisely the case in which a static formula has no basis
for preferring one channel over the other.

\paragraph{Quality dimension fusion.}
Three of the six quality dimensions are acoustically informed:
\texttt{professionalism} and \texttt{empathy} (agent channel),
and \texttt{customer\_satisfaction} (customer channel, final third of the
call weighted $2\times$ to emphasize how the call was left). For a channel
with mean valence $\bar v \in [-1, 1]$:
\begin{align}
  a &= 3 + 2\bar v \qquad \text{(maps $[-1,1] \to [1,5]$)} \\
  w_a &= W_{\text{cap}} \cdot \text{coverage}, \qquad W_{\text{cap}} = 0.4 \\
  \text{fused} &= \operatorname{clamp}\big(\operatorname{round}((1-w_a)\cdot \text{text} + w_a \cdot a),\ 1,\ 5\big)
\end{align}
where \emph{coverage} is the fraction of that channel's sentences the
sentiment model successfully processed. The text score is capped at a
maximum acoustic weight of 0.4, since the rubric's behaviours are defined
in text; acoustics modulate rather than override. When a call has no
acoustic data, $w_a = 0$ and the dimension is text-only.

\paragraph{Escalation risk fusion.}
The customer channel's escalation-score trajectory yields two summary
statistics: $\text{late\_mean}$ (mean score over the final third of the
call) and $\text{peak}$ (max score anywhere). An acoustic risk tier is
assigned from $\text{late\_mean}$ alone:
\[
  \text{acoustic\_risk} =
  \begin{cases}
    \text{escalate} & \text{late\_mean} \ge 0.50 \\
    \text{review}   & \text{late\_mean} \ge 0.30 \\
    \text{none}     & \text{otherwise}
  \end{cases}
\]
This is then merged against the LLM's text-derived risk tier:
\begin{itemize}
  \item If both tiers \textbf{agree}, the merge is deterministic
        (\texttt{method = agreement}).
  \item If they \textbf{disagree}, the call is routed to an
        \texttt{arbitrate} node -- a focused LLM call that is given both
        tier verdicts, the summary statistics, and the trajectory itself,
        and must return a single adjudicated risk level with a rationale
        (\texttt{method = llm\_arbitration}).
\end{itemize}
\texttt{peak} is retained only as contextual input to the arbitrator; it
is not used as a tier trigger (see \S\ref{subsec:formula} for why).

\subsection{Evolution of Compliance, Quality, and Workflow Evaluation}

\begin{itemize}
  \item \textbf{Compliance and quality} moved from a single monolithic
        prompt to independent, parallel LangGraph nodes, one per rubric
        section, each validated against its own Pydantic schema and
        retried independently on malformed output. This isolates failures:
        a malformed compliance response no longer invalidates the quality
        or workflow extraction for the same call.
  \item \textbf{Workflow checking} evaluates a per-call, LLM-generated
        checklist of expected steps (derived from a call subject
        extracted in a separate, isolated phase) against the transcript.
        Two correctness issues were found and fixed this cycle:
        \begin{enumerate}
          \item \emph{Subject leakage.} The subject-extraction prompt was
                allowed to surface call-specific figures (dollar amounts,
                account numbers, dates), which made generated checklist
                items trivially satisfiable by the same figures already
                in the transcript. The prompt now explicitly excludes
                call-specific detail from the extracted subject.
          \item \emph{False-vs-null asymmetry.} The step-checking prompt
                gave the model an escape hatch to mark a step
                \texttt{null} without matching guidance for when to
                affirmatively conclude \texttt{false}, which biased the
                model toward avoiding determinations. An explicit rule
                was added: \texttt{null} is reserved for steps the
                transcript genuinely cannot decide; a step that plainly
                never happened is \texttt{false}.
          \item \emph{Silent completeness failure.} Under upstream
                rate-limiting, a retried LLM call could return fewer
                checklist entries than steps given, and the prior
                index-based join silently padded the missing entries
                with \texttt{met = null} -- a truncated response still
                passed schema validation despite being incomplete. The
                workflow node now explicitly validates that the returned
                entry count matches the given step count and raises on
                mismatch, triggering the existing retry loop.
        \end{enumerate}
        After both fixes, a 22-call validation batch showed workflow
        pass rates spread meaningfully from 38\% to 100\% (previously an
        artificially high, largely null-contaminated result), with zero
        null entries remaining across the batch.
  \item \textbf{Escalation risk} gained a conditional arbitration step
        (\S\ref{subsec:formula}) in place of a fixed merge rule, changing
        the fusion step from a single formula for all cases to a
        two-path design: deterministic on agreement, model-adjudicated on
        disagreement.
\end{itemize}

\subsection{Current Fusion Formula / Logic}
\label{subsec:formula}

The escalation thresholds are empirically fitted rather than asserted.
Across a 22-call validation batch ($n=21$ calls with paired acoustic
data), the observed distribution of $\text{late\_mean}$ was: minimum
0.166, p25 0.219, median 0.227, p75 0.257, p90 0.281, maximum 0.406 (a
single outlier; the next-highest value was 0.281). The \texttt{review}
threshold (0.30) was set just above the calm cluster's p90, so only the
outlier call crosses it; the \texttt{escalate} threshold (0.50) was set
deliberately outside the observed range, since no call in the batch
warranted an acoustic-only escalate verdict and the bar should stay high
until data demonstrates otherwise.

\texttt{peak} was evaluated as a candidate second trigger but discarded:
its median across the batch was 0.558 and its p90 was 0.696, meaning
even calm, well-handled calls routinely spike past what would have been
its threshold. It has no discriminating power as a tier trigger and is
now recorded only as contextual input to the arbitrator.

Validating this design on the full batch: of the 22 calls, 18 fused by
deterministic agreement, 1 had no acoustic data (text-only), and 3 were
routed to arbitration. All 3 arbitrated cases resolved by favoring the
acoustic signal's read of the call's trajectory over the text-only
verdict -- including one case where a text-calm transcript was escalated
to \texttt{review} on the strength of a late-call acoustic escalation
spike, which then triggered the pipeline's investigation step.

\subsection{LLM Provider Evaluation}

Before wiring a fixed provider into the extraction pipeline, six
candidate provider/model pairs were benchmarked on the same structured
extraction task against two contrasting calls (one clean, high-scoring
banking call; one weaker banking call), so the choice of primary model
would be evidence-based rather than assumed. Each candidate's raw output
was auto-scored on four criteria:
\begin{enumerate}
  \item \textbf{Schema validity} -- does the output parse and validate
        against the Pydantic rubric schema?
  \item \textbf{Evidence fidelity} -- what fraction of quoted evidence
        strings are verbatim substrings of the transcript (vs.\
        hallucinated)?
  \item \textbf{Score spread} -- does the model's quality scoring
        actually discriminate between the clean and the weaker call, or
        does it return near-identical scores regardless of call quality?
  \item \textbf{Conservative bias} -- of the compliance items marked
        \texttt{PASS}, how many carry supporting evidence (a model that
        passes items without evidence is over-lenient)?
\end{enumerate}

\begin{table}[h]
\centering
\begin{tabular}{lccc}
\hline
\textbf{Provider / Model} & \textbf{Latency} & \textbf{Evidence Fidelity} & \textbf{Outcome} \\
\hline
Mistral (\texttt{mistral-small-latest})            & moderate & high, conservative & \textbf{Primary} -- only candidate whose scores discriminated between the two calls ([5,4,3,5,2] vs.\ [4,4,3,4,2]) \\
SambaNova (\texttt{Meta-Llama-3.3-70B-Instruct})    & fastest (2.6s) & 75\% & \textbf{Fallback} -- fastest, best raw fidelity, perfectly conservative \\
GitHub (\texttt{gpt-4o-mini})                       & moderate & reliable JSON & \textbf{Safety} -- always-on free baseline \\
SambaNova (\texttt{DeepSeek-V3.1})                  & moderate & 50\% & \textbf{Excluded} -- fabricates quotes, the worst failure mode for a compliance system \\
Gemini (\texttt{gemini-2.0-flash})                  & -- & -- & Benchmarked, not selected \\
Cohere (\texttt{command-r7b-12-2024})               & -- & -- & Benchmarked, not selected \\
\hline
\end{tabular}
\caption{LLM provider/model comparison for the QA extraction task
(\texttt{benchmark.py}). Ranking criteria: schema validity, verbatim
evidence fidelity, cross-call score spread, and conservative compliance
bias.}
\end{table}

The three retained tiers (Mistral $\to$ Llama-3.3-70B $\to$
GPT-4o-mini) were wired into a LiteLLM \texttt{Router} with strict
priority fallback: a request is served by the primary tier, and on a
429/timeout/error transparently drops to the next tier with per-tier
retry and cooldown, so a single provider outage or rate limit does not
fail the pipeline. This routing policy is derived directly from the
benchmark results above rather than chosen independently of them.

\subsection{Transcription: Provider Evaluation and Accuracy Metrics}

No structural changes were made to the transcription pipeline this
cycle. Work was limited to accuracy benchmarking of the current
transcription approach (per-channel faster-whisper \texttt{small.en},
stereo channel attribution) as a baseline for evaluating candidate LLM
transcription providers.

\begin{table}[h]
\centering
\begin{tabular}{lccc}
\hline
\textbf{Metric} & \textbf{Value} \\
\hline
Calls evaluated & 22 \\
Total reference words & 32{,}059 \\
Normalised WER & 9.8\% \\
Normalised accuracy & \textbf{90.2\%} \\
Target ($\geq 90\%$) & PASS \\
\hline
\end{tabular}
\caption{Overall transcription accuracy, normalised (Whisper
\texttt{EnglishTextNormalizer}: canonicalises numbers, contractions, and
spelling variants; strips disfluency annotations).}
\end{table}

\begin{table}[h]
\centering
\begin{tabular}{lccc}
\hline
\textbf{Domain} & \textbf{Calls} & \textbf{Accuracy (normalised)} & \textbf{Pass?} \\
\hline
Banking & 10 & 91.1\% & Yes \\
Health  & 7  & 86.5\% & No \\
Telecom & 5  & 93.2\% & Yes \\
\hline
\end{tabular}
\caption{Transcription accuracy by domain. Health falls short of the
90\% target, driven substantially by one outlier call
(\texttt{en\_CA\_Health\_1590851}, 74.4\% normalised accuracy, 37.9\%
customer-channel WER) rather than a uniform domain weakness.}
\end{table}
